% Study sheet: Volume I, Chapter 13, Model Serving.
% Numbers: \val{...} from model_serving_studysheet_numbers.tex, written by model_serving_studysheet_numbers.py.
\documentclass{studysheet}
\input{model_serving_studysheet_numbers}
\sheetvolume{1}
\sheetchapter{13}{Model Serving}
\sheetsources{Numbers: mlsysim and the chapter's modeled ResNet-50 profile on one H100; M/M/1 queueing; assumed inputs marked.}

\begin{document}
% ===================================================================== FRONT
\frontside
\begin{bigidea}
Training is judged by work finished per hour; serving is judged by whether \hl{each request meets its deadline}.
The model is one stage of the request path, live traffic decides when work arrives, and capacity is sized for the
\hl{slow tail}, not the average.
\end{bigidea}

\begin{leftcol}
\sheetsec{Key ideas}
\begin{ideas}
\idea{Serving inverts training.}{Requests arrive alone, so batch one exposes fixed costs a big batch shared,
  and the goal shifts from full use to headroom.}
\idea{The deadline is a budget.}{Network, decoding, queue, copies, model and response spend it in series;
  a faster model shrinks only its own slice.}
\idea{Utilization turns into waiting.}{Time in system grows as $1/(1-\ServingUtilization)$ and its p99 is
  \val{tailfactor}$\times$ the mean: size capacity at the utilization the tail permits.}
\idea{Batching trades waiting for throughput.}{A batch shares one weight read, but each request waits for it to
  form; at batch one, compiling beats lower precision.}
\idea{Replicas start late and can skew.}{Loading, compiling and warmup take seconds, so
  scale-out lags; preprocessing unlike training's silently changes answers.}
\idea{Language models.}{Prefill sets time to first token; decode rereads all weights per token,
  so bandwidth sets its pace; the key-value (KV) cache caps concurrency.}
\end{ideas}

\vspace{10pt}
\sheetsec{Where one request's time goes}
{\centering
\begin{tikzpicture}[sheet]
  \pgfmathsetmacro\u{18.5}
  \pgfmathsetmacro\xa{46}
  \pgfmathsetmacro\xb{\xa+\val{sjpeg}*\u}
  \pgfmathsetmacro\xc{\xb+\val{sresize}*\u}
  \pgfmathsetmacro\xd{\xc+\val{snorm}*\u}
  \pgfmathsetmacro\xe{\xd+\val{scopy}*\u}
  \pgfmathsetmacro\xf{\xe+\val{sfwd}*\u}
  \pgfmathsetmacro\xg{\xf+\val{stop}*\u}
  \pgfmathsetmacro\xff{\xe+\val{sfwdfast}*\u}
  \pgfmathsetmacro\xgg{\xff+\val{stop}*\u}
  % eager engine
  \node[lab, anchor=east] at (\xa-3,31) {eager};
  \draw[nbox] (\xa,24) rectangle (\xb,38);
  \draw[nbox] (\xb,24) rectangle (\xe,38);
  \draw[hotbox] (\xe,24) rectangle (\xf,38);
  \draw[nbox] (\xf,24) rectangle (\xg,38);
  \node[lab] at ({(\xa+\xb)/2},31) {JPEG decode \val{jpeg}};
  \node[lab, text=white] at ({(\xe+\xf)/2},31) {model \val{forward}};
  \node[note, anchor=west] at (\xg+3,31) {\val{total} ms};
  % model 10x faster
  \node[lab, anchor=east] at (\xa-3,9) {model \val{speedup}$\times$};
  \draw[nbox] (\xa,2) rectangle (\xb,16);
  \draw[nbox] (\xb,2) rectangle (\xe,16);
  \draw[hotbox] (\xe,2) rectangle (\xff,16);
  \draw[nbox] (\xff,2) rectangle (\xgg,16);
  \node[note, anchor=west] at (\xgg+3,9) {\val{totalfast} ms};
  \draw[dim] ({(\xb+\xe)/2},-1) -- ({(\xb+\xe)/2},-6);
  \node[small, anchor=north] at ({(\xb+\xe)/2+30},-7) {resize, normalize, copy to GPU: \val{restprep} ms};
  \draw[dim] (\xf+1,39) -- (\xf+1,44);
  \node[small, anchor=south] at (\xf,45) {top-5 \val{topfive}};
\end{tikzpicture}\par}
\vspace{6pt}
{\fontsize{8.6}{11}\selectfont Time each stage before optimizing. Idle gaps on the accelerator's timeline point to
the host; waiting that grows with load points to the queue.\par}
\end{leftcol}%
\begin{rightcol}
\fcard{Queueing: utilization against the tail (M/M/1)}%
  {\RequestTimeInSystem = \frac{\text{service time}}{1-\ServingUtilization},\qquad
   \TailRequestLatency \approx \val{tailfactor}\,\RequestTimeInSystem}%
  {Mean time in system and its p99 (the time 99\% of requests beat), with utilization $\ServingUtilization = \ArrivalRate/\ServiceRate$ (arrival rate over service rate).}%
  {To size: $\ServingUtilization \le 1 - \val{tailfactor}\times\text{service time}/\text{p99 budget}$; servers
   $=\ArrivalRate/(\ServingUtilization\,\ServiceRate)$, round up; apply headroom, round up again. Little's law, $\RequestConcurrency =
   \ArrivalRate\,\RequestTimeInSystem$, counts requests in flight. Assumes Poisson arrivals.}%
  {\begin{tikzpicture}[sheet]
     \pgfmathsetmacro\sx{210}
     \pgfmathsetmacro\sy{2}
     \draw[dim] (20,0) -- (20+\sx,0);
     \draw[dim] (20,0) -- (20,84);
     \foreach \t in {0,0.25,0.5,0.75,1}{
       \draw[dim] (20+\t*\sx,0) -- (20+\t*\sx,-2.5);
       \node[small, anchor=north] at (20+\t*\sx,-4) {\t};}
     \node[small, anchor=north] at (20+0.5*\sx,-13) {utilization $\ServingUtilization$};
     \node[small, anchor=south, rotate=90] at (16,42) {in service times};
     \draw[gridline] (20,{\val{cbudget}*\sy}) -- (20+\sx,{\val{cbudget}*\sy});
     \draw[accent, dashed, line width=0.6pt] ({20+\val{crho}*\sx},0) -- ({20+\val{crho}*\sx},78);
     \draw[muted, dashed, line width=0.6pt] plot[domain=0:0.975, samples=60] ({20+\x*\sx},{\sy/(1-\x)});
     \draw[accent, line width=0.9pt] plot[domain=0:0.885, samples=60] ({20+\x*\sx},{\val{tailfactor}*\sy/(1-\x)});
     \filldraw[fill=accent, draw=accent] ({20+\val{crho}*\sx},{\val{cbudget}*\sy}) circle (2.2pt);
     \node[note, anchor=south west] at (24,{\val{cbudget}*\sy+1}) {budget \val{wbudget} ms $=$ \val{wbudgetmult} service times};
     \node[note, anchor=north east] at ({20+\val{crho}*\sx-3},76) {max $\ServingUtilization$ \val{wrho}};
     \node[lab, text=accent, anchor=west] at ({20+0.89*\sx},78) {p99};
     \node[lab, text=muted, anchor=north west] at ({20+0.86*\sx},12) {mean};
   \end{tikzpicture}}

\vspace{12pt}
\sheetsec{Rules for batches, tokens and cost}
\begin{sheettable}{@{}l L@{}}
\thead{Quantity} & \thead{Rule (example)}\\ \midrule
Batch service time & $\BatchServiceTime(\BatchSize) \approx$ fixed $+$ per-item $\times\,\BatchSize$
  (ResNet-50, eager: \val{tfixed} $+$ \val{tper}$\,\BatchSize$ ms)\\
Batched capacity & $\EffectiveServiceRate(\BatchSize) = \BatchSize/\BatchServiceTime(\BatchSize)$; utilization
  $\ServingUtilization = \ArrivalRate/\EffectiveServiceRate(\BatchSize)$ must stay below one\\
Batched latency & mean $\BatchingWindow/2 + \BatchServiceTime(\BatchSize)$ (window $\BatchingWindow$, full
  batches, idle worker); the tail waits the full $\BatchingWindow$\\
Decode floor & $\DecodeTokenTimeFloor \approx \DataVolume/\AcceleratorMemoryBandwidth$,: weight bytes over memory bandwidth\\
KV cache per token & $2\times$ layers $\times$ KV heads $\times$ head width $\times$ bytes per value;
  conversations $\le$ (memory $-$ weights) $\div$ KV per conversation; one GiB $=$ \val{gibgb} GB\\
Cost per inference & (replicas $\times$ price per hour) $\div$ requests served per hour within the SLO\\
\end{sheettable}
\end{rightcol}

% ===================================================================== BACK
\backside
\begin{answers}
\ans{Finish}{\val{wpeak} $\div$ \val{trho} $\approx$ \val{tneed} req/s; one GPU does \val{tpergpu}/s, so \val{traw} $\to$ \val{tgpus}, and $\times\val{whead} = \val{theadraw} \to$ \val{tfinal} GPUs.
  A \val{tfaster}$\times$ faster service needs \val{tfewer}$\times$ fewer: it also raises the safe utilization.}
\ans{1}{Capacity $\BatchSize/\BatchServiceTime$: batch \val{bone}: \val{qcapone}/s (unstable); \val{beight}: \val{qcapeight}/s,
  $\ServingUtilization=\val{qrhoeight}$ (too hot); \val{bsixteen}: \val{qcapsixteen}/s, $\ServingUtilization=\val{qrhosixteen}$.
  Mean latency \val{qhalf} $+$ \val{qtsixteen} $=$ \val{qlat} ms; in flight \val{qrate} $\times$ \val{qlat} ms $\approx$
  \val{qinflight}.}
\columnbreak
\ans{2}{\val{lpertok} B $\times$ \val{lctx} $=$ \val{lconv} GB each; \val{lcapgib} GiB $=$ \val{lcapgb} GB, minus
  \val{lweights} leaves \val{lfree} GB: about \val{lconvs}. Floor \val{lweights} GB $\div$ \val{hbw} TB/s $\approx$
  \val{lfloor} ms (\val{ltokrate} tokens/s). Stutter is decode: bandwidth, not compute.}
\ans{3}{Per-stage timestamps or a timeline: accelerator gaps while the host decodes mean host-bound; wait before
  service that grows with load means queueing.}
\ans{4}{T4: \val{cload}/\val{ct4use} $=$ \val{ct4ratio} $\to$ \val{ct4n}, \$\val{ct4hour}/h, \$\val{ct4million} per
  million. H100: \val{chbusy}\% busy, \$\val{chmillion} per million (\$\val{chsat} only if full).}
\end{answers}

\begin{leftcol}
\sheetsec{Worked example}
\begin{practicebox}\boxtitle{How many accelerators serve the peak within the p99?}
The chapter's ResNet-50 plan: \val{wpeak} requests/s at peak, a \val{wbudget} ms p99 budget for queue plus
accelerator, \val{wservice} ms batch-one service time on one H100, \val{wheadpct}\% headroom.
\begin{steps}
\item Safe utilization: $1 - \val{tailfactor}\times\val{wservice}/\val{wbudget} = $ \res{\val{wrho}}.
\item Required rate: \val{wpeak} $\div$ \val{wrho} $\approx$ \res{\val{wneed} req/s}.
\item One GPU: one request per \val{wservice} ms $\approx$ \val{wpergpu} req/s, so \val{wneed} $\div$ \val{wpergpu} $=$
  \val{wraw} $\to$ \res{\val{wgpus} GPUs}.
\item Headroom: \val{wgpus} $\times$ \val{whead} $=$ \val{wheadraw} $\to$ \res{\val{wfinal} GPUs}.
\end{steps}
\takeaway{Size from the tail at a safe utilization, then add headroom; average throughput alone undersizes.}
\end{practicebox}

\vspace{12pt}
\sheetsec{Your turn: finish it}
\begin{practicebox}\boxtitle{What does a faster engine save?}
Same plan, but a graph executor removes Python dispatch: \val{tservice} ms per request.\par\vspace{4pt}
\textbf{Given:} safe utilization $1 - \val{tailfactor}\times\val{tservice}/\val{wbudget} = \val{trho}$.\par
\textbf{You:} find the required rate, the GPU count with \val{wheadpct}\% headroom, and why the fleet shrinks by
more than the service time did.
\end{practicebox}
\end{leftcol}%
\begin{rightcol}
\sheetsec{Common mistakes}
\mistake{``Make the model \val{speedup}$\times$ faster and requests get \val{speedup}$\times$ faster.''}%
  {Decoding, resizing, normalizing and copying (\val{hostshare}\% of the \val{total} ms request) are untouched: the request falls
   to \val{totalfast} ms, \val{amdahl}$\times$ faster.}{}
\mistake{``Run the accelerators at \val{utilhigh}\% to save money.''}%
  {Mean time in system goes from \val{meanlow}$\times$ (at \val{utillow}\%) to \val{meanhigh}$\times$ the service time;
   with \val{wservice} ms service, p99 goes from \val{ptailow} ms to \val{ptailhigh} ms.}{}
\mistake{``INT8 makes batch-one serving \val{fourx}$\times$ faster.''}%
  {Compiling the eager model (\val{rteager} ms) to one captured graph gives \val{rtfp32} ms, \val{rtcompile}$\times$;
   INT8 then adds \val{rtprecision}$\times$ (\val{rtint8} ms), because a per-kernel floor dominates one image.
   Precision pays at large batch.}{}


\vspace{2pt}
\sheetsec{Check yourself}
\begin{qs}
\item \assumed{} \val{qrate} requests/s reach one ResNet-50 replica with an \val{qwin} ms batching window and
  $\BatchServiceTime(\BatchSize) = \val{tfixed} + \val{tper}\,\BatchSize$ ms. Which batch cap (\val{bone}, \val{beight} or \val{bsixteen}) keeps
  utilization at or below \val{qmax}? With full batches and an idle worker, what is its mean latency, and how many requests are in flight?
\item Llama 3 8B (\val{llayers} layers, \val{lkvheads} KV heads of width \val{lheaddim}, 16-bit) holds
  \val{lweights} GB of weights on an \val{lcapgib} GiB H100 (\val{hbw} TB/s). How many \val{lctx}-token
  conversations fit? What is the batch-one floor per output token? A fast first token, then a stutter: which phase, and does
  more compute help?
\item A node misses its p99 while the accelerator is busy only a third of the time. What observation separates a
  slow host path from queueing?
\item \assumed{} \val{cload} images/s, kept at or below \val{cutil}\% utilization. T4s (\$\val{ct4price}/h,
  \val{ct4rate}/s each) or one H100 (\$\val{chprice}/h, \val{chrate}/s)? Cost per million images?
\end{qs}
\end{rightcol}
\end{document}
